\chapter{Linee guida e principi dell'interazione uomo-macchina}
\label{cap:hci}
\emergencystretch=2.5em

Un sistema può essere corretto, efficiente e ben progettato all'interno, e risultare comunque difficile da usare. Nel
corso degli anni sono state raccolte molte \textbf{linee guida} e molti \textbf{principi} per progettare buone interfacce
utente: nascono dall'esperienza, sono stati raffinati per decenni e valgono per la maggior parte dei sistemi
interattivi. Questo capitolo li presenta seguendo la lezione su \textit{Guidelines and Principles in Human-Computer
Interaction} (HCI, interazione uomo-macchina): le otto regole d'oro di Shneiderman e le dieci euristiche di Nielsen e
Molich, una per una, con gli esempi delle slide. Le finestre di esempio sono ridisegnate e tradotte in italiano.

% Finestre di esempio (mock-up)
% \begin{finestra}[larghezza]{titolo} contenuto \end{finestra}
\NewDocumentEnvironment{finestra}{O{\linewidth} m}{%
  \begin{tcolorbox}[enhanced, breakable=false, width=#1, nobeforeafter, box align=top, colback=white,
    colframe=swgray!70, boxrule=0.6pt, arc=1.5pt, left=5pt, right=5pt, top=4pt, bottom=4pt,
    fontupper=\scriptsize\Lato, colbacktitle=swteal, coltitle=white, fonttitle=\scriptsize\Lato,
    toptitle=1pt, bottomtitle=1pt, title={#2\hfill{\tiny --\enspace$\square$\enspace$\times$}}]%
  \raggedright}{\end{tcolorbox}}
% Pulsanti, campi e caselle delle finestre
\newcommand{\bottone}[2][swteal]{\tikz[baseline=(b.base)]\node[rectangle, rounded corners=1pt, fill=#1, text=white,
  font=\tiny\Lato\bfseries, inner sep=2.5pt, minimum width=1.2cm] (b) {#2};}
\newcommand{\bottonechiaro}[1]{\tikz[baseline=(b.base)]\node[rectangle, rounded corners=1pt, draw=swgray, fill=white,
  text=black!80, font=\tiny\Lato, inner sep=2.5pt, minimum width=1.2cm] (b) {#1};}
\newcommand{\pulsanti}[1]{\par\smallskip{\hfill #1\par}}
\newcommand{\campo}[2][2.2cm]{\tikz[baseline=(c.base)]\node[rectangle, draw=swgray!70, fill=white, font=\tiny\Lato,
  minimum width=#1, text width=#1-6pt, inner sep=2.5pt, align=left] (c) {#2};}
\newcommand{\casella}{\tikz[baseline=-0.4ex]\draw[line width=0.4pt, draw=swgray] (0,-0.03) rectangle (0.17,0.14);}
% Giudizio sotto una finestra
\newcommand{\male}{\tikz[baseline=(v.base)]\node[chip=swred] (v) {da evitare};}
\newcommand{\bene}{\tikz[baseline=(v.base)]\node[chip=swgreen] (v) {meglio};}
\newcommand{\benissimo}{\tikz[baseline=(v.base)]\node[chip=swgreen!70!black] (v) {ancora meglio};}

% ============================================================
\section{Linee guida e principi}

Le linee guida e i principi per progettare buone interfacce:

\begin{itemize}
  \item si applicano alla maggior parte dei sistemi interattivi;
  \item derivano dall'esperienza e sono stati raffinati nel corso di decenni;
  \item non hanno la pretesa di essere completi o universali: vanno verificati e adattati ai singoli domini di
        progettazione;
  \item restano comunque molto utili, per chi studia e per chi progetta.
\end{itemize}

Molti autori ne hanno proposto degli insiemi, più o meno noti:

\begin{itemize}
  \item le otto regole d'oro (\textit{Eight Golden Rules}) di Ben Shneiderman~[1];
  \item le dieci euristiche di usabilità (\textit{Ten Usability Heuristics}) di Jakob Nielsen e Rolf
        Molich~[2];
  \item i venti principi di usabilità studiati da Holcomb e Tharp~[3];
  \item gli otto principi di progettazione di Polson e Lewis, per interfacce in cui l'utente indovina che cosa
        fare~[4];
  \item le diciannove domande di Carroll e Rosson per l'analisi degli artefatti~[5].
\end{itemize}

Altre euristiche si trovano in~[6]. Questi insiemi servono in due momenti: per \textbf{guidare la
progettazione} e per \textbf{valutare} un'interfaccia esistente, cercandone i problemi di usabilità. Si sovrappongono in
parte, perché certi temi ricorrono in tutti (prevenire gli errori, ridurre il carico sulla memoria, \ldots). Qui ne
vediamo due, le euristiche di Nielsen e Molich e le regole d'oro di Shneiderman, e li percorriamo insieme, come
\textbf{unione} dei due insiemi.

\subsection{Le otto regole d'oro di Shneiderman}

Ben Shneiderman le ha formulate a partire dal 1985:

\begin{enumerate}
  \item cercare la \textbf{coerenza} (\textit{Strive for Consistency});
  \item cercare l'\textbf{usabilità universale} (\textit{Seek Universal Usability});
  \item offrire un \textbf{feedback informativo} (\textit{Offer Informative Feedback});
  \item progettare dialoghi che arrivino a una \textbf{conclusione} (\textit{Design Dialogs to Yield Closure});
  \item \textbf{prevenire gli errori} (\textit{Prevent Errors});
  \item permettere di \textbf{annullare} facilmente le azioni (\textit{Permit Easy Reversal of Actions});
  \item lasciare agli utenti il \textbf{controllo} (\textit{Keep Users in Control});
  \item ridurre il carico sulla \textbf{memoria a breve termine} (\textit{Reduce Short-term Memory Load}).
\end{enumerate}

\subsection{Le dieci euristiche di Nielsen e Molich}

Jakob Nielsen e Rolf Molich le hanno proposte nel 1990, e Nielsen le ha riviste nel 1994:

\begin{enumerate}
  \item dialogo \textbf{semplice e naturale} (\textit{Simple and Natural Dialogue});
  \item parlare la \textbf{lingua dell'utente} (\textit{Speak the User's Language});
  \item ridurre il carico sulla \textbf{memoria} dell'utente (\textit{Minimize User Memory Load});
  \item \textbf{coerenza} (\textit{Consistency});
  \item \textbf{feedback};
  \item \textbf{uscite} ben segnalate (\textit{Clearly Marked Exits});
  \item \textbf{scorciatoie} (\textit{Shortcuts});
  \item buoni \textbf{messaggi di errore} (\textit{Good Error Messages});
  \item \textbf{prevenire} gli errori (\textit{Prevent Errors});
  \item \textbf{aiuto} e documentazione (\textit{Help and Documentation}).
\end{enumerate}

\subsection{L'unione dei due elenchi}

La tabella mostra in quale regola e in quale euristica compare ciascun tema; le sezioni che seguono li trattano
nell'ordine della tabella.

\begin{center}
\renewcommand{\arraystretch}{1.25}
\begin{tabular}{@{}lcc@{}}
  \toprule
  \textbf{Tema} & \textbf{Shneiderman} & \textbf{Nielsen e Molich} \\ \midrule
  Dialogo semplice e naturale & -- & 1 \\
  Parlare la lingua dell'utente & -- & 2 \\
  Ridurre il carico sulla memoria & 8 & 3 \\
  Coerenza & 1 & 4 \\
  Feedback & 3 & 5 \\
  Uscite ben segnalate e azioni reversibili & 6, 7 & 6 \\
  Scorciatoie & -- & 7 \\
  Buoni messaggi di errore & -- & 8 \\
  Prevenire gli errori & 5 & 9 \\
  Aiuto e documentazione & -- & 10 \\
  Usabilità universale & 2 & -- \\
  Dialoghi che arrivano a una conclusione & 4 & -- \\
  \bottomrule
\end{tabular}
\end{center}

\info{Le euristiche di Nielsen nella versione del 1994}{
  Molti testi riportano le euristiche nella forma rivista da Nielsen nel 1994, con nomi diversi ma idee in gran parte
  uguali: visibilità dello stato del sistema; corrispondenza fra il sistema e il mondo reale; controllo e libertà
  dell'utente; coerenza e standard; prevenzione degli errori; riconoscere piuttosto che ricordare; flessibilità ed
  efficienza d'uso; design estetico e minimalista; aiutare gli utenti a riconoscere, diagnosticare e correggere gli
  errori; aiuto e documentazione. Chi le incontra in questa forma ritrova gli stessi temi di questo capitolo.
}

% ============================================================
\section{Dialogo semplice e naturale}

L'interfaccia dovrebbe essere \textbf{il più semplice possibile} (ma non più semplice di così!).

\begin{itemize}
  \item \textbf{Meno è meglio} (\textit{less is more}): ogni funzione o informazione in più è una cosa in più da
        imparare, da fraintendere e da scorrere quando si cerca quella che serve.
  \item Gli utenti alle prime armi possono essere sopraffatti da troppe informazioni.
  \item Lo dice anche la \textbf{legge di Hick}: il tempo che una persona impiega per prendere una decisione cresce in
        modo logaritmico con il numero e la complessità delle scelte.
\end{itemize}

\info{La legge di Hick in formula}{
  Nella forma più nota (legge di Hick-Hyman), con $n$ scelte equiprobabili il tempo di decisione è $T = a + b \log_2(n
  + 1)$, dove $a$ e $b$ sono costanti che si misurano sperimentalmente. Raddoppiare le voci di un menu non raddoppia il
  tempo di scelta, ma lo allunga sempre: ogni voce inutile costa qualcosa a chi deve scegliere.
}

L'interfaccia, poi, dovrebbe adattarsi nel modo più \textbf{naturale} possibile al compito dell'utente, con
corrispondenze (\textit{mappings}) e metafore: la versione digitale di un modulo è organizzata come quella di carta, e
un'app bussola funziona come una bussola vera.

\warning{{Meno è meglio, finché meno non diventa peggio}}{
  Una vignetta delle slide mostra il rovescio della medaglia: un'interfaccia «semplificata» riducendo tutto a un solo
  pulsante, che fa cose diverse se lo si tiene premuto per cinque secondi, se lo si tocca due volte, se lo si preme a
  metà\ldots{} La complessità non è sparita: si è spostata nella testa dell'utente. Lo stesso vale per un orologio
  intelligente con un solo tasto e troppe funzioni.
}

\subsection{Un dialogo poco naturale}

Sulla piattaforma \textit{docenti.unina} le notizie pubblicate da un docente sono ordinate dalla più recente e divise
in pagine. Il collegamento \textit{Risultati precedenti} porta però alle notizie \textbf{più recenti}, e \textit{Risultati
successivi} a quelle \textbf{più vecchie}: «precedente» e «successivo» seguono l'ordine delle pagine, non quello del
tempo, che è quello a cui pensa chi legge le notizie.

\begin{center}
\begin{tikzpicture}[
    notizia/.style={rectangle, draw=swgray!50, fill=swlight, font=\tiny\Lato, text width=5.6cm, inner sep=3pt,
                    align=left}]
  \node[notizia] (n1) at (0,0) {\textbf{Avviso esame del 12 febbraio} \hfill pubblicato il 30/01};
  \node[notizia, below=1pt of n1] (n2) {\textbf{Spostamento della lezione di giovedì} \hfill pubblicato il 21/01};
  \node[notizia, below=1pt of n2] (n3) {\textbf{Risultati della prova intercorso} \hfill pubblicato il 13/01};
  \node[font=\tiny\Lato, text=swblue, below left=4pt and 0pt of n3.south east] (succ) {\underline{Risultati successivi}};
  \node[font=\tiny\Lato, text=swblue, left=0.9cm of succ] (prec) {\underline{Risultati precedenti}};
  \draw[-{Stealth[length=1.8mm]}, swred, line width=0.6pt] (prec.south) -- ++(0,-0.35)
    node[below, font=\tiny, text=swred, align=center] {notizie\\più recenti};
  \draw[-{Stealth[length=1.8mm]}, swred, line width=0.6pt] (succ.south) -- ++(0,-0.35)
    node[below, font=\tiny, text=swred, align=center] {notizie\\più vecchie};
\end{tikzpicture}
\end{center}

% ============================================================
\section{Parlare la lingua dell'utente}

Il dialogo dovrebbe usare parole, frasi e concetti \textbf{familiari all'utente}, invece di termini orientati al
sistema: è la progettazione centrata sulle persone (\textit{human-centered design}).

\begin{itemize}
  \item I dialoghi dovrebbero essere nella lingua madre dell'utente (\textbf{localizzazione}): e non solo il testo, ma
        anche gli elementi non verbali, come le icone.
  \item Bisogna fare attenzione alle parole che si usano. Progettando per il grande pubblico si usano parole che tutti
        capiscono, nel loro significato comune; progettando per un gruppo di utenti con una propria terminologia
        specialistica, si usano proprio i loro termini tecnici.
\end{itemize}

Nell'esempio di sinistra il calcolo del prezzo di un noleggio parla la lingua di chi ha scritto il programma: il
«cliente» (ma l'utente pensa a sé come a chi noleggia), una sigla per esperti come CDW, un «preventivo» invece di un
prezzo. A destra la sigla è sciolta e spiegata:

\Needspace{16\baselineskip}
\begin{center}
\begin{minipage}[t]{0.38\linewidth}\centering
  \begin{finestra}{Calcolo del prezzo del noleggio}
    {\small\textbf{Calcola il prezzo del noleggio}}\par\smallskip
    Età del cliente: \campo[1.2cm]{18}\par\smallskip
    Modello: \campo[1.6cm]{Modello \hfill$\blacktriangledown$}\par\smallskip
    \casella\ CDW
    \pulsanti{\bottone{Calcola preventivo}}
  \end{finestra}\par\smallskip\male
\end{minipage}\hspace{0.06\linewidth}%
\begin{minipage}[t]{0.42\linewidth}\centering
  \begin{finestra}{Calcolo del prezzo del noleggio}
    {\small\textbf{Calcola il prezzo del noleggio}}\par\smallskip
    Età di chi noleggia: \campo[1.2cm]{18}\par\smallskip
    Modello: \campo[1.6cm]{Modello \hfill$\blacktriangledown$}\par\smallskip
    \casella\ Copertura dei danni da collisione (CDW)\par
    {\tiny Seleziona questa opzione per acquistare la copertura dei danni da collisione e non dover pagare gli
    eventuali danni all'auto noleggiata.}
    \pulsanti{\bottone{Calcola il prezzo}}
  \end{finestra}\par\smallskip\bene
\end{minipage}
\end{center}

Parlare la lingua dell'utente vuol dire anche guardare le interazioni \textbf{dal suo punto di vista}: chi ha appena
comprato delle azioni non pensa che qualcuno gliele abbia vendute.

\Needspace{9\baselineskip}
\begin{center}
\begin{minipage}[t]{0.4\linewidth}\centering
  \begin{finestra}{Operazione riuscita}
    Ti abbiamo venduto 100 azioni di XYZ Corp.
    \pulsanti{\bottonechiaro{OK}}
  \end{finestra}\par\smallskip\male
\end{minipage}\hspace{0.06\linewidth}%
\begin{minipage}[t]{0.4\linewidth}\centering
  \begin{finestra}{Operazione riuscita}
    Hai comprato 100 azioni di XYZ Corp.
    \pulsanti{\bottonechiaro{OK}}
  \end{finestra}\par\smallskip\bene
\end{minipage}
\end{center}

% ============================================================
\section{Ridurre il carico sulla memoria}

I computer ricordano le cose con grande precisione; la memoria di lavoro delle persone, molto meno. L'interfaccia
dovrebbe togliere all'utente il peso di ricordare, il più possibile. Come?

\begin{itemize}
  \item \textbf{Riconoscere è meglio che ricordare}: mostrare le possibilità tra cui scegliere, invece di chiedere di
        richiamarle a memoria.
  \item Quando si chiede un dato, descriverne il \textbf{formato} e dare un \textbf{esempio}; dichiarare esplicitamente
        gli intervalli e i valori ammessi, se ce ne sono.
\end{itemize}

Nel primo esempio un sistema comunica un numero di conferma di diciotto cifre e, alla schermata dopo, chiede di
inserirlo: l'utente non dovrebbe doverlo ricordare (né ricopiarlo), perché il sistema lo conosce già.

\Needspace{12\baselineskip}
\begin{center}
\begin{minipage}[t]{0.36\linewidth}\centering
  \begin{finestra}{Rinnovo}
    \centering
    {\color{swred}\textbf{La tua richiesta di rinnovo\\è stata approvata!}}\par\smallskip
    Il tuo numero di conferma è:\par
    \textbf{168426246998723652}\par\smallskip
    Ora puoi procedere con il pagamento.
    \pulsanti{\bottone{Vai al pagamento}}
  \end{finestra}
\end{minipage}\hspace{0.04\linewidth}\raisebox{-1.2cm}{\Large\textcolor{swteal}{$\Rightarrow$}}\hspace{0.04\linewidth}%
\begin{minipage}[t]{0.36\linewidth}\centering
  \begin{finestra}{Fatturazione e pagamenti}
    Inserisci il numero di conferma dell'acquisto per procedere:\par\smallskip
    \campo[3.6cm]{\phantom{0}}
    \pulsanti{\bottone{Passo successivo >}}
  \end{finestra}\par\smallskip\male
\end{minipage}
\end{center}

Nel secondo si deve indicare una delle province italiane con la sua sigla. Un campo di testo obbliga a ricordare il
codice (qual è quello di Lecce? E quello di Lecco?); un elenco a discesa con i nomi per esteso e le sigle toglie
all'utente il peso di ricordare i codici di tutte le province.

\Needspace{12\baselineskip}
\begin{center}
\begin{minipage}[t]{0.36\linewidth}\centering
  \begin{finestra}{Prodotto}
    Provincia di origine del prodotto (sigla):\par\smallskip
    \campo[1.6cm]{MI}
  \end{finestra}\par\smallskip\male
\end{minipage}\hspace{0.08\linewidth}%
\begin{minipage}[t]{0.36\linewidth}\centering
  \begin{finestra}{Prodotto}
    Provincia di origine del prodotto (sigla):\par\smallskip
    \campo[2.9cm]{\hfill$\blacktriangledown$}\par
    \tikz\node[rectangle, draw=swgray!70, fill=white, font=\tiny\Lato, text width=2.9cm-6pt, inner sep=2.5pt,
      align=left] {ANCONA (AN)\\AOSTA (AO)\\AREZZO (AR)\\ASCOLI PICENO (AP)\\ASTI (AT)\\AVELLINO (AV)\\\ldots};
  \end{finestra}\par\smallskip\bene
\end{minipage}
\end{center}

% ============================================================
\section{Coerenza}

Per essere usabile un sistema deve essere coerente, sia all'interno sia con il mondo esterno.

\begin{itemize}
  \item \textbf{Coerenza interna}, all'interno del prodotto o di una famiglia di prodotti:
        \begin{itemize}
          \item situazioni simili dovrebbero richiedere sequenze di azioni simili: eliminare un cliente ed eliminare un
                fornitore dovrebbero richiedere gli stessi passi;
          \item la stessa informazione dovrebbe essere presentata nello stesso modo e nello stesso punto in tutte le
                schermate.
        \end{itemize}
  \item \textbf{Coerenza esterna}, con le convenzioni ormai consolidate.
\end{itemize}

Alcuni esempi dalle slide:

\begin{itemize}
  \item su \textit{docenti.unina} il controllo \textit{Risultati successivi} cambia posizione: nella prima pagina è il
        primo da sinistra, nelle pagine seguenti si sposta a destra, perché davanti compare \textit{Risultati
        precedenti}; chi lo cerca nello stesso punto clicca quello sbagliato;
\end{itemize}

\begin{center}
\begin{tikzpicture}[
    barra/.style={rectangle, draw=swgray!50, fill=swlight, minimum width=6.2cm, minimum height=0.5cm},
    link/.style={font=\tiny\Lato, text=swblue}]
  \node[barra] (b1) at (0,0) {};
  \node[font=\tiny\Lato\bfseries, anchor=east] at ($(b1.west)+(-0.1,0)$) {Pagina 1};
  \node[link, anchor=west] (x1) at ($(b1.west)+(0.15,0)$) {\underline{Risultati successivi}};
  \node[barra] (b2) at (0,-0.8) {};
  \node[font=\tiny\Lato\bfseries, anchor=east] at ($(b2.west)+(-0.1,0)$) {Pagine seguenti};
  \node[link, anchor=west] (p2) at ($(b2.west)+(0.15,0)$) {\underline{Risultati precedenti}};
  \node[link, anchor=west] (x2) at ($(p2.east)+(0.25,0)$) {\underline{Risultati successivi}};
  \draw[-{Stealth[length=1.8mm]}, swred, dashed, line width=0.6pt] (x1.south) to[bend right=10] (x2.north);
\end{tikzpicture}
\end{center}

\begin{itemize}
  \item nella procedura di accettazione degli appuntamenti di un servizio sanitario americano (dal sito del Nielsen
        Norman Group, \url{https://www.nngroup.com/articles/consistency-and-standards/}) le schermate successive
        cambiano disposizione e aspetto dei pulsanti, e l'utente deve ricominciare ogni volta a cercarli;
  \item un buon esempio di coerenza interna è la suite Microsoft Office: Word, Excel e PowerPoint hanno la stessa barra
        multifunzione, con le stesse schede e gli stessi comandi negli stessi punti, e chi conosce un programma si
        orienta subito negli altri;
  \item un esempio di coerenza esterna sono i siti di commercio elettronico: la casella di ricerca in alto, il logo a
        sinistra, il carrello in alto a destra. Chi ha usato un negozio online sa già dove cercare le cose in uno nuovo.
\end{itemize}

% ============================================================
\section{Feedback}

Il sistema dovrebbe informare \textbf{continuamente} l'utente su che cosa sta facendo e su come sta interpretando i suoi
comandi:

\begin{itemize}
  \item non solo quando si verificano errori: un feedback positivo è importante quanto uno negativo;
  \item quando è possibile, il feedback va dato anche se il sistema si guasta;
  \item il peggior feedback possibile è l'assenza di feedback.
\end{itemize}

Il feedback, poi, non dovrebbe essere troppo astratto e generico. I tre avvisi che seguono riguardano la stessa
copia di file; solo l'ultimo dà all'utente tutto ciò che gli serve per decidere:

\Needspace{14\baselineskip}
\begin{center}
\begin{minipage}[t]{0.27\linewidth}\centering
  \begin{finestra}{File duplicato}
    La copia sovrascriverà un file esistente.
    \pulsanti{\bottone{Continua}\ \bottone[swred]{Annulla}}
  \end{finestra}\par\smallskip\male
\end{minipage}\hfill
\begin{minipage}[t]{0.27\linewidth}\centering
  \begin{finestra}{File duplicato}
    La copia sovrascriverà il file «softeng.txt».
    \pulsanti{\bottone{Continua}\ \bottone[swred]{Annulla}}
  \end{finestra}\par\smallskip\bene
\end{minipage}\hfill
\begin{minipage}[t]{0.38\linewidth}\centering
  \begin{finestra}{File duplicato}
    La cartella di destinazione contiene già un file «softeng.txt».\par\smallskip
    Vuoi sostituire il file esistente\par
    \hspace*{1em}{\tiny Dimensione: 350 KB \quad Modificato: 1 set 2024, 17:45}\par
    con questo?\par
    \hspace*{1em}{\tiny Dimensione: 580 KB \quad Modificato: 12 set 2024, 11:37}
    \pulsanti{\bottone{Continua}\ \bottone[swred]{Annulla}}
  \end{finestra}\par\smallskip\benissimo
\end{minipage}
\end{center}

\subsection{Quanto deve durare un feedback}

Feedback diversi richiedono una diversa \textbf{persistenza}:

\begin{itemize}
  \item alcuni contano solo per la durata di un fenomeno, o sono semplici conferme che un'operazione è stata eseguita:
        possono sparire da soli, come i messaggi \textit{toast};
  \item altri, soprattutto avvisi ed errori, richiedono che l'utente li confermi esplicitamente;
  \item altri ancora devono essere molto persistenti, e fare parte stabilmente dell'interfaccia.
\end{itemize}

\Needspace{10\baselineskip}
\begin{center}
\tcbset{persist/.style={nobeforeafter, colback=white, colframe=swgray!60, boxrule=0.5pt, arc=1.5pt, left=3pt,
  right=3pt, top=2pt, bottom=2pt, fontupper=\tiny\Lato}}
\begin{minipage}[t]{0.2\linewidth}\centering
  \tcbox[persist, colback=swgreen!80, colframe=swgreen!80, fontupper=\tiny\Lato\color{white}]{La password è stata
    aggiornata}\par\smallskip
  {\scriptsize\bfseries sparisce da sola}
\end{minipage}\hfill
\begin{minipage}[t]{0.25\linewidth}\centering
  \begin{tcolorbox}[persist]
    Download del file. Tempo rimanente: 00:35\par
    \tikz{\fill[swlight] (0,0) rectangle (2.5,0.12); \fill[swgreen] (0,0) rectangle (1.33,0.12);}\hfill 53\%
  \end{tcolorbox}\par\smallskip
  {\scriptsize\bfseries mostra l'avanzamento}
\end{minipage}\hfill
\begin{minipage}[t]{0.27\linewidth}\centering
  \begin{tcolorbox}[persist]
    \textcolor{q5}{\textbf{!}}\ Non è stato possibile applicare alcune impostazioni per un conflitto con i criteri di
    sicurezza.\par{\hfill\bottonechiaro{OK}\par}
  \end{tcolorbox}\par\smallskip
  {\scriptsize\bfseries va confermato}
\end{minipage}\hfill
\begin{minipage}[t]{0.22\linewidth}\centering
  \begin{tcolorbox}[persist, halign=center]
    Stato del remoto Git: \textcolor{swred}{\textbf{online}}\par\smallskip
    \bottone{Pull}\ \bottone{Push}
  \end{tcolorbox}\par\smallskip
  {\scriptsize\bfseries sempre visibile}
\end{minipage}
\end{center}

\subsection{Feedback e tempi di risposta}

Il feedback è cruciale quando il sistema impiega molto a rispondere. I limiti da tenere a mente sono tre:

\begin{center}
\renewcommand{\arraystretch}{1.3}
\begin{tabular}{@{}P{2.6cm}P{11.4cm}@{}}
  \toprule
  \textbf{Tempo} & \textbf{Che cosa percepisce l'utente, e che feedback serve} \\ \midrule
  meno di 0,1 s & la reazione sembra istantanea: basta mostrare il risultato o confermare l'esito \\
  meno di 1 s & il filo del pensiero non si interrompe, anche se la reazione non sembra più istantanea: non serve un
    feedback speciale \\
  circa 10 s & è il limite per tenere l'attenzione dell'utente sul compito: oltre i 10 secondi il feedback è
    indispensabile. Conviene stimare quando l'operazione finirà (l'utente vorrà fare altro intanto) e aggiornare spesso
    l'indicatore di avanzamento \\
  \bottomrule
\end{tabular}
\end{center}

\subsection{Il bias della percezione del lavoro}

Il \textbf{bias della percezione del lavoro} (\textit{labor perception bias}) dice che le persone si fidano di più, e
danno più valore, alle cose di cui percepiscono il lavoro che c'è dietro. A nessuno piace aspettare; ma se l'utente ha
aspettative alte (perché si tratta di denaro, del backup o della migrazione di dati importanti, di un'analisi o di un
rapporto) un'attesa troppo breve lo rende sospettoso.

\begin{itemize}
  \item Una \textbf{schermata di lavoro} (\textit{labor screen}) subito dopo un'azione importante, che mostra che cosa
        il sistema sta facendo, può migliorare l'esperienza d'uso.
  \item A volte i progettisti aggiungono degli «inganni benevoli», come tempi di caricamento finti: ne parla un articolo
        di \textit{The Atlantic} del 2017,
        \url{https://www.theatlantic.com/technology/archive/2017/02/why-some-apps-use-fake-progress-bars/517233/}.
\end{itemize}

\Needspace{17\baselineskip}
\begin{center}
\tcbset{schermo/.style={nobeforeafter, box align=top, width=\linewidth, height=4.6cm, colback=white, colframe=black!70,
  boxrule=1pt, arc=4pt, left=5pt, right=5pt, top=5pt, bottom=5pt, fontupper=\tiny\Lato}}
\begin{minipage}[t]{0.25\linewidth}
  \begin{tcolorbox}[schermo]
    \textbf{Cerca voli}\par
    Trova i voli migliori per la tua destinazione!\par\medskip
    Data:\par\campo[2.6cm]{\phantom{0}}\par
    Aeroporto di partenza:\par\campo[2.6cm]{\phantom{0}}\par
    Aeroporto di arrivo:\par\campo[2.6cm]{\phantom{0}}\par\medskip
    {\centering\bottone[swblue]{Cerca i voli migliori}\par}
  \end{tcolorbox}
\end{minipage}\hspace{0.02\linewidth}\raisebox{-2.2cm}{\Large\textcolor{swteal}{$\Rightarrow$}}\hspace{0.02\linewidth}%
\begin{minipage}[t]{0.25\linewidth}
  \begin{tcolorbox}[schermo, colback=swblue!85, fontupper=\tiny\Lato\color{white}]
    \vspace*{1.1cm}
    {\centering\textbf{Cerco i voli migliori\ldots}\par}\medskip
    \textbullet\ Contatto le compagnie\par
    \textbullet\ Contatto le agenzie\par
    \textbullet\ Tratto le offerte\par
    \textbullet\ Scelgo i voli migliori
  \end{tcolorbox}
\end{minipage}\hspace{0.02\linewidth}\raisebox{-2.2cm}{\Large\textcolor{swteal}{$\Rightarrow$}}\hspace{0.02\linewidth}%
\begin{minipage}[t]{0.25\linewidth}
  \begin{tcolorbox}[schermo]
    \textbf{Ecco i voli migliori per te!}\par\medskip
    1. EasyJet (23 €)\par\hspace*{0.8em}diretto, parte alle 6:30\par\smallskip
    2. Ryanair (31 €)\par\hspace*{0.8em}diretto, parte alle 8:30\par\smallskip
    3. Wizz Air (55 €)\par\hspace*{0.8em}diretto, parte alle 14:00
  \end{tcolorbox}
\end{minipage}
\par\smallskip
{\scriptsize\hspace*{0.3\linewidth}\textcolor{swteal}{\textbf{la schermata di lavoro, in mezzo, mostra che cosa sta
facendo il sistema}}}
\end{center}

Senza la schermata centrale i risultati comparirebbero all'istante: proprio per questo l'utente potrebbe dubitare che
siano davvero i migliori. Con la schermata di lavoro vede il sistema contattare compagnie e agenzie, e si fida di più del
risultato.

% ============================================================
\section{Uscite ben segnalate e azioni reversibili}

Gli utenti vogliono sentirsi \textbf{in controllo} dell'interazione, e commetteranno comunque degli errori, per quanto
il sistema sia ben progettato. Il sistema dovrebbe quindi offrire un modo semplice per \textbf{uscire} dalla maggior
parte delle situazioni:

\begin{itemize}
  \item se un'azione non si può completare entro 10 secondi, l'utente dovrebbe poterla interrompere e annullare;
  \item nelle operazioni con effetti collaterali, l'uscita si offre con una funzione di \textbf{annullamento}
        (\textit{Undo}), che riporta il sistema allo stato precedente.
\end{itemize}

\Needspace{9\baselineskip}
\begin{center}
\begin{minipage}[t]{0.42\linewidth}\centering
  \begin{finestra}{Addestramento del modello}
    Addestramento avviato. Tempo rimanente: circa 125 ore.\par\smallskip
    \tikz{\fill[swlight] (0,0) rectangle (4.6,0.12); \fill[swgreen] (0,0) rectangle (0.3,0.12);}
    \pulsanti{\bottone[swred]{Annulla}}
  \end{finestra}
\end{minipage}\hspace{0.06\linewidth}%
\begin{minipage}[t]{0.36\linewidth}\centering
  \begin{finestra}{Contatti eliminati}
    Hai eliminato 3653 contatti.
    \pulsanti{\bottone{OK}\ \bottone[swred]{Annulla l'operazione}}
  \end{finestra}
\end{minipage}
\end{center}

Un articolo citato nelle slide (\url{https://www.uxmatters.com/mt/archives/2020/03/are-you-sure-versus-undo-design-and-technology.php})
confronta le due strategie: chiedere prima «Sei sicuro?» oppure eseguire subito e permettere di annullare. Torneremo
sulla questione parlando degli errori: per gli errori più comuni la seconda funziona meglio.

% ============================================================
\section{Scorciatoie}

In generale, usare un'interfaccia dovrebbe richiedere la conoscenza di poche regole. Gli utenti esperti, però,
dovrebbero poter eseguire le azioni frequenti con \textbf{scorciatoie} e \textbf{acceleratori}:

\begin{itemize}
  \item tasti funzione o combinazioni di tasti che racchiudono un intero comando in una sola pressione;
  \item il doppio clic su un oggetto per eseguire l'azione più comune su quell'oggetto;
  \item pulsanti specifici per raggiungere le funzioni importanti proprio nei punti del dialogo in cui servono più
        spesso;
  \item il riuso della storia dell'interazione, per ripetere rapidamente gli stessi comandi;
  \item buoni valori predefiniti nei moduli, quando è possibile.
\end{itemize}

Le slide mostrano tre esempi dagli ambienti di sviluppo: in IntelliJ IDEA il menu dei \textbf{progetti recenti} (un
acceleratore basato sulla storia) e le combinazioni di tasti mostrate accanto alle voci dei menu contestuali, così che
l'utente le impari usandole; in VS Code la scheda di riferimento delle scorciatoie da tastiera.

\subsection{Scorciatoie che si imparano}

Le scorciatoie da tastiera dovrebbero essere facili da \textbf{imparare} e da \textbf{ricordare}: non basta scegliere
combinazioni di tasti qualsiasi, purché diverse. Se l'utente deve consultare il riferimento ogni volta che vuole usare
una scorciatoia, quella non è una scorciatoia. In IntelliJ IDEA le operazioni di refactoring \textit{Introduce \ldots}
hanno combinazioni costruite sull'iniziale di ciò che si introduce:

\begin{center}
\renewcommand{\arraystretch}{1.2}
\begin{tabular}{@{}ll@{}}
  \toprule
  \textbf{Operazione} & \textbf{Scorciatoia} \\ \midrule
  Introduce \textbf{V}ariable & Ctrl+Alt+\textbf{V} \\
  Introduce \textbf{C}onstant & Ctrl+Alt+\textbf{C} \\
  Introduce \textbf{F}ield & Ctrl+Alt+\textbf{F} \\
  Introduce \textbf{P}arameter & Ctrl+Alt+\textbf{P} \\
  \bottomrule
\end{tabular}
\end{center}

\subsection{Scorciatoie coerenti}

Le scorciatoie devono anche essere \textbf{coerenti}. Ctrl+X, Ctrl+C e Ctrl+V sono interpretati universalmente come
taglia, copia e incolla fin dai tempi dell'Apple Lisa (1983). La ricerca, invece, è Ctrl+F in molti programmi (VS Code,
per esempio) ma non in tutti: le slide citano Microsoft Office, dove la ricerca è Ctrl+Maiusc+T, ed Excel usato dentro
Teams, dove è Ctrl+Alt+T. Chi passa da un programma all'altro sbaglia.

% ============================================================
\section{Buoni messaggi di errore}

I messaggi di errore sono critici per l'usabilità:

\begin{itemize}
  \item rappresentano situazioni in cui l'utente è in difficoltà e potrebbe non riuscire a usare il sistema per
        raggiungere i suoi obiettivi;
  \item sono un'occasione per aiutarlo a capire meglio il sistema: di solito chi incontra un errore è molto motivato a
        leggere il messaggio.
\end{itemize}

\dfn{Le quattro regole di Shneiderman per i messaggi di errore}{
  \begin{enumerate}
    \item Devono essere formulati in un \textbf{linguaggio chiaro}, evitando codici oscuri.
    \item Devono essere \textbf{precisi}, non vaghi o generici.
    \item Devono aiutare l'utente a \textbf{risolvere} il problema in modo costruttivo.
    \item Devono essere \textbf{cortesi}, senza intimidire o colpevolizzare l'utente.
  \end{enumerate}
}

\subsection{Linguaggio chiaro, messaggi precisi}

Il codice di un errore SQL, o il nome tecnico del vincolo violato, non dicono nulla a chi sta registrando un cliente;
nemmeno un generico «non è stato possibile salvare» lo aiuta. Il messaggio di destra dice che cosa non è andato, perché
e che cosa fare:

\Needspace{16\baselineskip}
\begin{center}
\begin{minipage}[t]{0.27\linewidth}\centering
  \begin{finestra}{Errore}
    SQLError 23503
    \pulsanti{\bottone{Continua}}
  \end{finestra}\par\smallskip
  \begin{finestra}{Errore}
    \textbf{FOREIGN KEY VIOLATION} when inserting new record in customers relation.
    \pulsanti{\bottone{Continua}}
  \end{finestra}\par\smallskip\male
\end{minipage}\hfill
\begin{minipage}[t]{0.27\linewidth}\centering
  \begin{finestra}{Errore}
    Non è stato possibile salvare il nuovo cliente.
    \pulsanti{\bottone{Continua}}
  \end{finestra}\par\smallskip\male
\end{minipage}\hfill
\begin{minipage}[t]{0.36\linewidth}\centering
  \begin{finestra}{Errore}
    Non è stato possibile salvare il nuovo cliente.\par\smallskip
    La città indicata («Npoli») non è un nome di città valido. Controlla di averlo scritto correttamente e
    riprova.\par\smallskip
    Se il problema persiste, contatta l'amministratore del sistema.
    \pulsanti{\bottone{Continua}}
  \end{finestra}\par\smallskip\bene
\end{minipage}
\end{center}

\subsection{Messaggi costruttivi e cortesi}

Un messaggio costruttivo non si limita a dire che c'è un problema, ma spiega come risolverlo:

\Needspace{12\baselineskip}
\begin{center}
\begin{minipage}[t]{0.36\linewidth}\centering
  \begin{finestra}{Errore}
    Impossibile creare il nuovo account: password non valida.
    \pulsanti{\bottone{Continua}}
  \end{finestra}\par\smallskip\male
\end{minipage}\hspace{0.06\linewidth}%
\begin{minipage}[t]{0.4\linewidth}\centering
  \begin{finestra}{Errore}
    Per motivi di sicurezza e di conformità, una password deve contenere almeno 16 caratteri. La password che hai
    inserito ne ha 7.\par\smallskip
    Inserisci una password più lunga e riprova.
    \pulsanti{\bottone{Continua}}
  \end{finestra}\par\smallskip\bene
\end{minipage}
\end{center}

Un messaggio cortese non intimidisce e non colpevolizza: chi non riesce a raggiungere il suo obiettivo si sente già
abbastanza a disagio. Vanno evitati i termini aggressivi, come «AZIONE ILLEGALE DELL'UTENTE!», «LAVORO INTERROTTO»,
«PROCESSO UCCISO», «ERRORE FATALE»; ed è meglio formulare il messaggio in modo da suggerire che il problema sia del
sistema. In un certo senso lo è davvero, perché un'interfaccia ben progettata avrebbe potuto prevenire l'errore.

\Needspace{9\baselineskip}
\begin{center}
\begin{minipage}[t]{0.42\linewidth}
\begin{lstlisting}[style=plain, name=Interprete]
$> Foo()
Errore fatale! Stai cercando di eseguire
una funzione (Foo) che non hai definito!
$>
\end{lstlisting}
\centering\male
\end{minipage}\hspace{0.06\linewidth}%
\begin{minipage}[t]{0.42\linewidth}
\begin{lstlisting}[style=plain, name=Interprete]
$> Foo()
Non so come eseguire Foo.
Sei sicuro di aver definito la funzione Foo?
$>
\end{lstlisting}
\centering\bene
\end{minipage}
\end{center}

\subsection{Messaggi a più livelli}

I messaggi di errore servono agli utenti, ma anche al personale tecnico che gestisce il sistema. Gli utenti in genere non
capiscono i dettagli tecnici (codici di errore, \textit{stack trace}), che al personale tecnico servono per individuare
il problema: un messaggio può doverli contenere entrambi.

\Needspace{14\baselineskip}
\begin{center}
\begin{minipage}{0.62\linewidth}
  \begin{finestra}{Errore}
    Si è verificato un errore durante la conversione dell'unità di misura. Il valore indicato («35°C») non è un numero
    decimale valido.\par\smallskip
    Codice di errore: \textbf{E457X}\par
    Stack trace:\par
    {\color{swred}\ttfamily\tiny Exception in thread "main" it.unina.InputEx\\
    \hspace*{1em}at it.unina.Scanner.throwF(Scan.java:53)\\
    \hspace*{1em}at it.unina.Scanner.next(Scan.java:112)\\
    \hspace*{1em}at it.unina.Scanner.nextF(Scan.java:152)\\
    \hspace*{1em}at it.unina.Util.validF(Util.java:156)\\
    \hspace*{1em}at it.unina.Util.isValid(Util.java:47)\\
    \hspace*{1em}at it.unina.Converter.getInp(Conv.java:6)}
    \pulsanti{\bottone{Continua}}
  \end{finestra}
\end{minipage}
\end{center}

Spesso è preferibile \textbf{separare} i livelli: l'utente comune non si spaventa per messaggi dall'aspetto strano, e il
personale tecnico trova comunque le informazioni che gli servono. La finestra di errore può anche contenere un
collegamento a un sito di assistenza.

\Needspace{12\baselineskip}
\begin{center}
\begin{minipage}[t]{0.44\linewidth}\centering
  \begin{finestra}{Errore}
    Si è verificato un errore durante la conversione dell'unità di misura. Il valore indicato («35°C») non è un numero
    decimale valido.
    \pulsanti{\bottonechiaro{Mostra i dettagli}\ \bottone{Continua}}
  \end{finestra}
\end{minipage}\hspace{0.03\linewidth}\raisebox{-0.9cm}{\Large\textcolor{swteal}{$\Rightarrow$}}\hspace{0.03\linewidth}%
\begin{minipage}[t]{0.4\linewidth}\centering
  \begin{finestra}{Dettagli dell'errore}
    Codice di errore: \textbf{E457X}\par
    {\color{swred}\ttfamily\tiny Exception in thread "main" it.unina.InputEx\\
    \hspace*{1em}at it.unina.Scanner.throwF(Scan.java:53)\\
    \hspace*{1em}at it.unina.Scanner.next(Scan.java:112)\\
    \hspace*{1em}at it.unina.Scanner.nextF(Scan.java:152)}
    \pulsanti{\bottone{Chiudi}}
  \end{finestra}
\end{minipage}
\end{center}

Lo fa, a modo suo, anche la celebre «schermata blu» di Windows, che le slide riportano con il codice di arresto
\texttt{SOFTWARE\_ENGINEERING\_LECTURE}: una frase semplice per l'utente («il tuo PC ha riscontrato un problema e deve
essere riavviato»), e un codice e un QR code per chi dovrà capire che cosa è successo.

% ============================================================
\section{Prevenire gli errori}

Ancora meglio di un buon messaggio di errore è \textbf{evitare} l'errore. Bisogna cercare di non mettere l'utente in
situazioni che favoriscono gli errori:

\begin{itemize}
  \item se gli si chiede di scrivere il nome di una città, c'è il rischio di errori di ortografia;
  \item se deve inserire un intervallo di date nel futuro, in un formato preciso, c'è il rischio che sbagli il formato o
        che inserisca una data già passata.
\end{itemize}

L'interfaccia va progettata per evitare (o ridurre al minimo) questi errori: per esempio con un calendario da cui
scegliere le date, come quello di Booking.com, che non permette di scegliere date passate. Non è solo una questione di
usabilità: di solito vuol dire anche meno lavoro per formalizzare i casi d'uso e meno codice per gestire le situazioni
di errore.

\subsection{Slip e mistake}

Secondo Don Norman esistono due categorie di errori:

\begin{itemize}
  \item gli \textbf{slip} (sviste): l'utente intende compiere un'azione, ma finisce per compierne un'altra. Preme il
        tasto «Invio» invece di «Backspace», o clicca sul pulsante «Riduci a icona» invece che su «Ingrandisci»;
  \item i \textbf{mistake} (sbagli): l'utente si pone un obiettivo inadatto al problema o al compito. Il responsabile di
        un sito di commercio elettronico vuole eliminare tutti gli articoli di una categoria, e crede che eliminando la
        categoria si eliminino anche i suoi articoli; in realtà gli articoli vengono spostati, senza dirlo, nella
        categoria «Altro».
\end{itemize}

Se l'utente si pone l'obiettivo giusto ma sbaglia l'esecuzione, ha fatto uno \textbf{slip}. Gli slip:

\begin{itemize}
  \item nascono di solito da comportamenti automatici;
  \item sono più frequenti negli utenti esperti (chi sta ancora imparando a usare un sistema fa più attenzione);
  \item sono spesso \textbf{errori di cattura} (\textit{capture errors}): quando due sequenze di azioni hanno un inizio
        comune e una delle due si usa molto più spesso dell'altra, l'utente si ritrova a seguire senza accorgersene
        quella più frequente, anche se voleva fare l'altra.
\end{itemize}

Per questo, in genere, è meglio permettere di \textbf{annullare} facilmente le azioni che affidarsi solo a una finestra
di conferma: chi elimina quasi sempre la cartella giusta preme «Continua» per abitudine, senza leggere.

\Needspace{11\baselineskip}
\begin{center}
\begin{minipage}[t]{0.42\linewidth}\centering
  \begin{finestra}{Conferma eliminazione}
    Vuoi davvero eliminare la cartella «SoftEng slides» e tutti i file che contiene?\par\smallskip
    L'operazione non può essere annullata.
    \pulsanti{\bottone{Continua}\ \bottone[swred]{Annulla}}
  \end{finestra}\par\smallskip\male
\end{minipage}\hspace{0.06\linewidth}%
\begin{minipage}[t]{0.42\linewidth}\centering
  \begin{tcolorbox}[enhanced, nobeforeafter, box align=top, width=\linewidth, colback=swred!85!black, colframe=swred!85!black,
    arc=3pt, left=6pt, right=6pt, top=5pt, bottom=5pt, fontupper=\scriptsize\Lato\color{white}]
    La cartella \textbf{«SoftEng slides»} e i \textbf{20 file} che conteneva sono stati eliminati.
    \par\smallskip{\hfill\bottonechiaro{Annulla l'operazione}\par}
  \end{tcolorbox}\par\smallskip\bene
\end{minipage}
\end{center}

I \textbf{mistake} sono molto più critici:

\begin{itemize}
  \item nascono spesso da un \textbf{modello mentale} del sistema sbagliato;
  \item possono essere molto più difficili da individuare, e quindi più pericolosi. Nell'esempio del sito di commercio
        elettronico, quando si accorgerà il responsabile che gli articoli sono ancora in vendita, nella categoria
        «Altro»?
\end{itemize}

% ============================================================
\section{Aiuto e documentazione}

Idealmente un sistema dovrebbe essere così facile da usare da non richiedere né aiuto né documentazione in aggiunta
all'interfaccia. Purtroppo non è sempre possibile: a parte i sistemi da usare al volo (\textit{walk-up-and-use}), la
maggior parte delle interfacce ha abbastanza funzioni da giustificare un manuale e, forse, un sistema di aiuto. Un
manuale può servire anche agli utenti abituali per acquisire una competenza maggiore e diventare più produttivi.

Avere un bel manuale e un buon sistema di aiuto, però, \textbf{non riduce} i requisiti di usabilità: «è tutto spiegato
nel manuale!» non è una buona scusa per un'interfaccia progettata male.

\dfn{La verità fondamentale sui manuali d'uso}{
  Gli utenti \textbf{non leggono} i manuali. Preferiscono passare il tempo in attività che li fanno sentire produttivi,
  e di solito cominciano a usare il sistema senza aver letto le istruzioni.

  \textbf{Corollario}: se un utente vuole leggere il manuale, probabilmente è nel panico e ha bisogno di aiuto
  immediato. In questi casi sono particolarmente utili i manuali in linea, con la consultazione orientata ai compiti
  e la ricerca.
}

% ============================================================
\section{Usabilità universale}

L'usabilità va cercata \textbf{per tutti}. Bisogna tenere presenti:

\begin{itemize}
  \item le differenze fra utenti \textbf{alle prime armi} ed \textbf{esperti};
  \item le diverse fasce di \textbf{età};
  \item le \textbf{disabilità};
  \item le differenze \textbf{internazionali}.
\end{itemize}

Progettare per tutti non significa ottenere un prodotto complessivamente meno efficace: spesso molte categorie di utenti
traggono vantaggio da accorgimenti pensati per le esigenze di una categoria specifica. È l'\textbf{effetto delle rampe
sui marciapiedi} (\textit{curb cut effect}): gli scivoli nati per le carrozzine servono anche a chi spinge un passeggino,
a chi trascina una valigia, ai ciclisti.

\begin{center}
\begin{tikzpicture}
  \node[align=center, font=\small\itshape, text width=11cm] (c) {«Fai un programma che anche un idiota può usare, e
    soltanto un idiota vorrà usarlo.»};
  \node[font=\scriptsize, anchor=north east] at (c.south east) {Arthur Bloch, \textit{Principio di Shaw}, \textit{La
    legge di Murphy}, 1977};
  \draw[swgray!60, line width=1.5pt] ($(c.north west)+(-0.1,0.05)$) -- ($(c.south west)+(-0.1,-0.35)$);
\end{tikzpicture}
\end{center}

La citazione, riportata nelle slide, ricorda il rischio opposto: semplificare tutto per il principiante più sprovveduto
fino a rendere il prodotto inutile per chiunque altro. Usabilità universale vuol dire servire bene sia i principianti
sia gli esperti (per esempio con le scorciatoie).

% ============================================================
\section{Dialoghi che arrivano a una conclusione}

Le sequenze di azioni vanno organizzate in gruppi con un \textbf{inizio}, una parte \textbf{centrale} e una
\textbf{fine}. Un feedback informativo al termine di ogni gruppo dà all'utente la soddisfazione di aver completato
qualcosa, un senso di sollievo e il segnale che può prepararsi al gruppo di azioni successivo. I siti di commercio
elettronico, per esempio, accompagnano il cliente attraverso una serie di passi chiari:

\begin{center}
\begin{tikzpicture}[
    passo/.style={rectangle, rounded corners=3pt, fill=#1, text=white, font=\scriptsize\bfseries, align=center,
                  minimum height=0.9cm, minimum width=2.7cm, inner sep=3pt},
    passo/.default=swteal]
  \node[passo] (p1) at (0,0) {1. Articoli\\nel carrello};
  \node[passo] (p2) at (3.4,0) {2. Pagamento\\e indirizzo};
  \node[passo] (p3) at (6.8,0) {3. Pagamento};
  \node[passo=swgreen] (p4) at (10.2,0) {Ordine confermato!};
  \foreach \a/\b in {p1/p2, p2/p3, p3/p4} { \draw[flow] (\a.east) -- (\b.west); }
\end{tikzpicture}
\end{center}

Il passo finale, con la conferma dell'ordine, chiude il dialogo: il cliente sa che l'acquisto è concluso e che può
dedicarsi ad altro.

% ============================================================
\section{Letture e riferimenti}

\begin{enumerate}[label={[\arabic*]}, leftmargin=2.2em, itemsep=2pt]
  \item B. Shneiderman, C. Plaisant, \textit{Designing the User Interface: Strategies for Effective
    Human-Computer Interaction}, Pearson Education, 2010.
  \item R. Molich, J. Nielsen, «Improving a human-computer dialogue», \textit{Communications of the ACM},
    33(3), 1990, pp. 338--348. \url{https://dl.acm.org/doi/10.1145/77481.77486}
  \item R. Holcomb, A. L. Tharp, «What users say about software usability», \textit{International
    Journal of Human-Computer Interaction}, 3(1), 1991, pp. 49--78. \url{https://doi.org/10.1080/10447319109525996}
  \item P. G. Polson, C. H. Lewis, «Theory-based design for easily learned interfaces»,
    \textit{Human-Computer Interaction}, 5(2--3), 1990, pp. 191--220. \url{https://doi.org/10.1080/07370024.1990.9667154}
  \item J. M. Carroll, M. B. Rosson, «Getting around the task-artifact cycle: how to make claims and
    design by scenario», \textit{ACM Transactions on Information Systems}, 10(2), 1992, pp. 181--212.
    \url{https://dl.acm.org/doi/abs/10.1145/146802.146834}
  \item J. Nielsen, «Enhancing the explanatory power of usability heuristics», \textit{Proceedings of
    the SIGCHI Conference on Human Factors in Computing Systems}, 1994, pp. 152--158.
    \url{https://dl.acm.org/doi/10.1145/191666.191729}
\end{enumerate}

\gbox{In sintesi}{azure-gradient-6!80!black}{
  \begin{itemize}
    \item Le linee guida di HCI nascono dall'esperienza: servono a progettare e a valutare le interfacce. Le più note
          sono le otto regole d'oro di \textbf{Shneiderman} e le dieci euristiche di \textbf{Nielsen e Molich}, che in
          gran parte si sovrappongono.
    \item \textbf{Dialogo semplice e naturale} (meno è meglio, la legge di Hick, ma senza esagerare) e \textbf{lingua
          dell'utente} (le sue parole, il suo punto di vista).
    \item \textbf{Memoria}: riconoscere è meglio che ricordare. \textbf{Coerenza} interna ed esterna.
    \item \textbf{Feedback} continuo, concreto, con la persistenza giusta; indispensabile oltre i 10 secondi (limiti di
          0,1, 1 e 10 secondi); a volte un'attesa visibile aumenta la fiducia (bias della percezione del lavoro).
    \item \textbf{Uscite e annullamento} per tenere l'utente in controllo; \textbf{scorciatoie} per gli esperti, facili da
          ricordare e coerenti.
    \item \textbf{Messaggi di errore} chiari, precisi, costruttivi e cortesi, magari a più livelli; meglio ancora
          \textbf{prevenire} gli errori. Gli \textbf{slip} si curano con l'annullamento, i \textbf{mistake} nascono da
          un modello mentale sbagliato.
    \item Gli utenti non leggono i manuali; l'usabilità va cercata per tutti; i dialoghi devono arrivare a una
          conclusione.
  \end{itemize}
}
